\documentclass[
    12pt,
    oneside,
    a4paper,
    brazil,
    article
]{abntex2}

\usepackage[T1]{fontenc}
\usepackage[utf8]{inputenc}
\usepackage{helvet}
\renewcommand{\familydefault}{\sfdefault}
\usepackage{babel}
\usepackage{booktabs}
\usepackage{array}
\usepackage{amsmath}
\usepackage{hyperref}
\usepackage{microtype}

\setlrmarginsandblock{3cm}{2cm}{*}
\setulmarginsandblock{3cm}{2cm}{*}
\checkandfixthelayout

\OnehalfSpacing
\setlength{\parindent}{1.25cm}
\setlength{\parskip}{0cm}

\hypersetup{
    colorlinks=true,
    linkcolor=black,
    urlcolor=black,
    citecolor=black,
    pdfauthor={Diogo Tadeu, Thiago Henrique, Wendel Cauã, Marcelo},
    pdftitle={Análise de Algoritmos para o Problema da Mochila 0-1 com Peso e Volume}
}

\titulo{Análise de Algoritmos para o Problema da Mochila 0-1 com Peso e Volume}
\autor{Diogo Tadeu\\Thiago Henrique\\Wendel Cauã\\Marcelo}
\local{Brasil}
\data{2026}
\instituicao{Trabalho Prático de PAA}
\tipotrabalho{Relatório}
\preambulo{Relatório apresentado como parte da avaliação da disciplina de Projeto e Análise de Algoritmos.}

\begin{document}
\selectlanguage{brazil}

\pretextual

\imprimircapa
\imprimirfolhaderosto*

\begin{resumo}
Este relatório apresenta a implementação e avaliação de três estratégias para o problema da Mochila 0-1 com duas restrições: peso e volume. Foram implementados Backtracking com podas, Programação Dinâmica com estados esparsos e Branch and Bound com limites superiores. Os experimentos avaliaram o tempo de execução em 20.113 execuções, distribuídas em 867 combinações de quantidade de itens, capacidade de peso e capacidade de volume. Os resultados mostram que a Programação Dinâmica foi a alternativa mais estável nas instâncias analisadas, enquanto o Backtracking apresentou os maiores picos de tempo. O Branch and Bound reduziu bastante o custo em relação ao Backtracking, mas ainda apresentou crescimento elevado em instâncias difíceis.

\vspace{\onelineskip}
\noindent
\textbf{Palavras-chave}: Mochila 0-1. Programação Dinâmica. Backtracking. Branch and Bound. Avaliação experimental.
\end{resumo}

\pdfbookmark[0]{\contentsname}{toc}
\tableofcontents*
\cleardoublepage

\textual

\section{Introdução}

O problema da Mochila 0-1 consiste em selecionar um subconjunto de itens de forma a maximizar o lucro total, respeitando restrições de capacidade. Na versão estudada neste trabalho, cada item possui três atributos inteiros positivos: lucro, peso e volume. A mochila possui uma capacidade máxima de peso $W$ e uma capacidade máxima de volume $V$. Como a escolha é 0-1, cada item pode ser escolhido no máximo uma vez.

O objetivo do trabalho foi implementar e comparar três algoritmos exatos para resolver o problema: Backtracking, Programação Dinâmica e Branch and Bound. A comparação foi feita principalmente pela métrica de tempo de execução, considerando diferentes tamanhos de entrada e diferentes capacidades de peso e volume. Além disso, foram produzidos arquivos CSV, análises estatísticas, gráficos e dashboards para apoiar a interpretação dos resultados.

De forma resumida, os resultados indicaram que a Programação Dinâmica apresentou maior estabilidade, principalmente por reutilizar estados já calculados. O Backtracking foi fortemente afetado pelo aumento do número de itens. O Branch and Bound teve desempenho intermediário: melhor que o Backtracking em muitos casos, mas ainda sensível à qualidade das podas. O restante do relatório está organizado em descrição dos algoritmos, análise de complexidade, avaliação experimental, conclusão e referências.

\section{Descrição dos algoritmos e complexidades}

\subsection{Backtracking com poda}

O Backtracking percorre uma árvore de decisão em que, para cada item, existem duas possibilidades: incluir ou não incluir o item na mochila. A implementação ordena os itens por uma densidade normalizada, considerando peso e volume relativos, e usa duas podas principais: descarte de ramos que excedem peso ou volume e descarte de ramos cujo lucro atual somado ao lucro restante possível não supera a melhor solução encontrada.

No pior caso, as podas podem não eliminar ramos suficientes, e o algoritmo ainda precisa analisar um número exponencial de subconjuntos. Assim, a complexidade de tempo no pior caso é:

\[
T(n) = O(2^n)
\]

em que $n$ é a quantidade de itens. A complexidade de espaço é:

\[
M(n) = O(n)
\]

considerando a pilha recursiva e o vetor de itens escolhidos no caminho atual.

\subsection{Programação Dinâmica com estados esparsos}

A Programação Dinâmica utilizada é iterativa, não recursiva com memoization. O algoritmo mantém estados identificados pelo par $(peso, volume)$, armazenando o maior lucro encontrado para cada combinação alcançável. Em vez de construir uma matriz tridimensional completa $dp[i][w][v]$, a implementação usa uma tabela hash, por meio de \texttt{unordered\_map}, para guardar apenas estados alcançados. Também são removidos estados dominados, isto é, estados que possuem peso e volume não melhores e lucro menor ou igual a outro estado.

A formulação clássica da mochila 0-1 com duas restrições possui tempo:

\[
T(n,W,V) = O(nWV)
\]

e espaço:

\[
M(W,V) = O(WV)
\]

quando se usa otimização por camadas, ou $O(nWV)$ se toda a matriz tridimensional for armazenada. Na implementação esparsa, o custo prático depende da quantidade $S$ de estados mantidos. Assim, o tempo observado pode ser descrito como $O(nS)$, mais o custo das podas de dominância, e o espaço como $O(S)$. No pior caso, $S$ pode se aproximar de $W \cdot V$, mantendo a natureza pseudo-polinomial do método.

\subsection{Branch and Bound}

O Branch and Bound também explora uma árvore de decisão 0-1, mas usa limites superiores para evitar a exploração de ramos que não podem melhorar a melhor solução conhecida. A implementação começa com uma solução inicial gulosa, o que aumenta o poder das podas desde o início. O limite superior combina três estimativas: valor restante máximo, limite fracionário por peso e limite fracionário por volume. A solução final continua sendo 0-1; a ideia fracionária é usada apenas para estimar um limite superior.

No pior caso, o Branch and Bound ainda pode visitar uma quantidade exponencial de nós. Como o cálculo do limite pode custar tempo proporcional ao número de itens restantes, a complexidade de tempo no pior caso é:

\[
T(n) = O(n \cdot 2^n)
\]

Na prática, as podas costumam reduzir bastante o número de nós visitados. A complexidade de espaço é $O(n)$, considerando a profundidade da recursão e os vetores de itens escolhidos, além das listas auxiliares ordenadas.

\section{Avaliação experimental}

\subsection{Configuração dos experimentos}

Os programas foram implementados em C++ e compilados por Make para execução no WSL. Foram gerados três executáveis separados, um para cada algoritmo, além de um gerador automático de instâncias e um executor de experimentos. As instâncias possuem itens com valores inteiros positivos, e os índices dos itens começam em 1. Para cada combinação solicitada de $n$, $W$ e $V$, foram geradas repetições aleatórias, permitindo calcular médias e desvios.

A saída dos experimentos foi registrada em CSV com, entre outros campos, quantidade de itens, capacidades, repetição, algoritmo, lucro obtido, itens escolhidos, tempo em segundos e status. Foi adotado limite de execução de 7200 segundos por algoritmo, permitindo interromper casos que ultrapassassem duas horas e prosseguir com o restante dos testes.

\subsection{Métrica de avaliação}

A principal métrica de avaliação foi o tempo de execução em segundos. Também foram registrados lucro máximo e itens escolhidos para verificar a consistência das soluções. Como os algoritmos são exatos, espera-se que, quando finalizados corretamente, retornem o mesmo lucro ótimo para a mesma instância. Para comparação estatística, foram aplicados testes por combinação de $n$, $W$ e $V$. Quando havia três algoritmos, foi usado o teste de Friedman; nas comparações par a par, foi usado Wilcoxon. Considerou-se diferença estatisticamente significativa quando $p < 0{,}05$ e empate estatístico quando $p \geq 0{,}05$.

\subsection{Resultados obtidos}

Ao todo, foram analisadas 20.113 execuções em 867 combinações de parâmetros. A Tabela~\ref{tab:resumo} resume os tempos médios e máximos por conjunto de experimentos e algoritmo.

\begin{table}[h]
\centering
\caption{Resumo global dos tempos de execução.}
\label{tab:resumo}
\small
\begin{tabular}{lrrrr}
\toprule
\textbf{CSV} & \textbf{Alg.} & \textbf{Exec.} & \textbf{Média (s)} & \textbf{Máx. (s)} \\
\midrule
100a600 & DP & 1403 & 0,25 & 5,27 \\
100a600 & BT & 1401 & 59,94 & 2686,18 \\
100a600 & B\&B & 1401 & 5,59 & 306,02 \\
\midrule
300 itens + W/V & DP & 1080 & 1,14 & 76,72 \\
300 itens + W/V & BT & 1080 & 0,80 & 31,85 \\
300 itens + W/V & B\&B & 1080 & 0,11 & 3,66 \\
\midrule
B\&B vs DP 1000 & DP & 5900 & 0,59 & 9,85 \\
B\&B vs DP 1000 & B\&B & 5900 & 11,40 & 1362,08 \\
\midrule
variação total & DP & 291 & 0,018 & 0,20 \\
variação total & BT & 288 & 260,06 & 13565,98 \\
variação total & B\&B & 289 & 17,81 & 631,14 \\
\bottomrule
\end{tabular}
\end{table}

Os resultados mostram grande diferença entre os algoritmos. No conjunto 100a600, por exemplo, o Backtracking teve média de 59,94 s, enquanto o Branch and Bound teve 5,59 s e a Programação Dinâmica 0,25 s. No conjunto variação total, o Backtracking atingiu média de 260,06 s e máximo superior a 13.565 s, evidenciando o crescimento exponencial em instâncias difíceis.

A Programação Dinâmica apresentou comportamento mais previsível e tempos menores em grande parte dos testes. Isso ocorreu porque o método reutiliza estados e evita enumerar subconjuntos completos. O Branch and Bound reduziu substancialmente o custo em relação ao Backtracking, mas ainda teve picos altos quando os limites superiores não foram suficientemente fortes para eliminar muitos ramos. Em alguns conjuntos com capacidades muito grandes, a Programação Dinâmica sentiu mais o aumento da quantidade de estados possíveis, o que explica casos em que B\&B foi competitivo ou mais rápido.

Nos testes estatísticos, a maioria das comparações indicou diferença significativa de tempo entre os algoritmos. As comparações DP vs BT apresentaram 249 diferenças significativas e 26 empates estatísticos; DP vs B\&B apresentou 733 diferenças e 132 empates; BT vs B\&B apresentou 274 diferenças e apenas 1 empate. O teste de Friedman, aplicado quando havia os três algoritmos, indicou diferença significativa em 275 combinações analisadas.

\section{Discussão}

O comportamento observado está de acordo com as complexidades teóricas. O Backtracking possui pior caso exponencial, e seus tempos cresceram muito quando $n$ aumentou. As podas simples ajudam, mas não eliminam a explosão combinatória. O Branch and Bound também possui pior caso exponencial, porém os limites superiores permitem cortar subárvores inteiras, explicando sua melhora em relação ao Backtracking. Ainda assim, quando muitos ramos parecem promissores, o algoritmo continua explorando uma quantidade elevada de nós.

A Programação Dinâmica tem comportamento pseudo-polinomial e depende das capacidades $W$ e $V$ e da quantidade de estados mantidos. A versão esparsa implementada reduziu o consumo de memória e o tempo em relação a uma matriz tridimensional completa, especialmente quando poucas combinações de peso e volume permaneceram relevantes. Por isso, nos testes com instâncias grandes, a DP foi geralmente a opção mais estável.

\section{Conclusão}

O trabalho implementou e comparou três algoritmos exatos para a Mochila 0-1 com restrições de peso e volume. A avaliação experimental mostrou que a escolha do algoritmo depende fortemente das características da instância. Para os testes realizados, a Programação Dinâmica com estados esparsos foi a estratégia mais estável e apresentou os menores tempos em grande parte dos casos. O Branch and Bound foi uma alternativa importante por reduzir bastante o custo do Backtracking, mas ainda apresentou picos de tempo em instâncias difíceis. O Backtracking, apesar de simples e útil para compreender a estrutura combinatória do problema, foi o menos adequado para instâncias grandes.

Como trabalhos futuros, seria possível medir também o número de nós visitados e o número de estados mantidos, além de comparar diferentes estratégias de ordenação, limites superiores e podas de dominância. Também seria interessante avaliar o consumo de memória de forma sistemática, já que esse fator é especialmente relevante para a Programação Dinâmica.

\begin{referencias}

\bibitem{cormen2009}
CORMEN, T. H.; LEISERSON, C. E.; RIVEST, R. L.; STEIN, C. \textit{Introduction to Algorithms}. 3. ed. Cambridge: MIT Press, 2009.

\bibitem{martello1990}
MARTELLO, S.; TOTH, P. \textit{Knapsack Problems: Algorithms and Computer Implementations}. Chichester: John Wiley \& Sons, 1990.

\bibitem{kellerer2004}
KELLERER, H.; PFERSCHY, U.; PISINGER, D. \textit{Knapsack Problems}. Berlin: Springer, 2004.

\bibitem{sipser2012}
SIPSER, M. \textit{Introduction to the Theory of Computation}. 3. ed. Boston: Cengage Learning, 2012.

\end{referencias}

\end{document}
